\documentclass[aspectratio=169,10pt]{beamer}
\usetheme{metropolis}
\usepackage{booktabs}
\title{Q01: From Next-Token Prediction to Capabilities}
\subtitle{What the objective can explain---and what it cannot}
\author{DLS 2026}
\date{October 2026}
\begin{document}
\maketitle
\begin{frame}{The question}
\begin{block}{Student question}
How does next-token prediction lead to problem solving, coding, and instruction following?
\end{block}
\begin{itemize}
\item Start with the training objective, not a metaphor.
\item Separate a plausible mechanism from a demonstrated claim.
\item End with an evaluation design.
\end{itemize}
\end{frame}
\begin{frame}{The objective is narrow}
\[\mathcal{L}(\theta)=-\sum_t \log p_\theta(x_{t+1}\mid x_{\leq t})\]
\begin{itemize}
\item Predict the next token from prior context.
\item No explicit ``reasoning,'' ``coding,'' or ``instruction following'' loss term.
\item A Transformer is one architecture for computing this conditional distribution.
\end{itemize}
\end{frame}
\begin{frame}{Why prediction rewards structure}
\begin{enumerate}
\item Text and code contain reusable regularities.
\item Better representations of these regularities reduce prediction loss.
\item Prompts supply examples and constraints at inference time.
\end{enumerate}
\alert{Important:} usefulness for prediction is not identical to human-like understanding.
\end{frame}
\begin{frame}{In-context learning is a concrete clue}
\begin{itemize}
\item A prompt can change behavior without changing weights.
\item Few-shot prompting demonstrated substantial task adaptation in large LMs.
\item Induction-head evidence links a particular circuit family to one in-context behavior in controlled models.
\end{itemize}
\vfill\small Brown et al. (2020); Olsson et al. (2022)
\end{frame}
\begin{frame}{Training data matters}
\begin{columns}
\column{.48\textwidth}
\textbf{Pre-training}
\begin{itemize}\item text, code, examples\item predicts continuations\item learns broad statistical structure\end{itemize}
\column{.48\textwidth}
\textbf{Post-training}
\begin{itemize}\item preference/instruction data\item shifts favored responses\item does not erase the distinction from pre-training\end{itemize}
\end{columns}
\end{frame}
\begin{frame}{Do not overclaim}
\begin{itemize}
\item A persuasive answer may be wrong.
\item A correct answer may be memorized, pattern-completed, or generalized.
\item ``Emergent'' can depend on the metric and threshold chosen.
\end{itemize}
\alert{Capability claims need held-out, contamination-aware evaluation.}
\end{frame}
\begin{frame}{Mini experiment}
Sequence: \texttt{A B A B A \ldots}
\begin{itemize}
\item Unigram predictor: choose the globally frequent token.
\item Context-copy predictor: find an earlier matching context and copy its continuation.
\item Compare accuracy on repeated-pattern sequences and on a changed rule.
\end{itemize}
\vfill The toy heuristic illustrates a context-dependent algorithm; it is not a model of all Transformer behavior.
\end{frame}
\begin{frame}{Takeaway}
\begin{block}{Best current answer}
Next-token training can reward representations and context-sensitive computations that make structured text, code, and demonstrations predictable. It provides a route to broad behavior, not a proof of robust reasoning or human-like intelligence.
\end{block}
\vfill\small References: Vaswani et al. (2017); Brown et al. (2020); Olsson et al. (2022); Wei et al. (2022); Schaeffer et al. (2023).
\end{frame}
\end{document}
